\newif\ifglobal

%%%%%%%%%%%%%%%%%%%%%%%
%%% Compile to view %%%
%%%%%%%%%%%%%%%%%%%%%%%

%%%%%%%%%%%%%%%%%%%%%%%%%%%%%%%%%%%%%%%%%%%%%%%%%%%
%%% Readme:                                     %%%
%%% https://www.overleaf.com/read/bwdjvfjksvjy/ %%%
%%%%%%%%%%%%%%%%%%%%%%%%%%%%%%%%%%%%%%%%%%%%%%%%%%%

% >>> M?: marks a module setting number or important notice

% >>> M4: Set {./relative-path} to external files for this module <<<
\def \modulefiles {./42-MIPIDSI}
% To add an external file, use:
% (Replace '---' with actual filename)
% '\input{\modulefiles/tables/---.tex}' for tables
% '\includegraphics{\modulefiles/figures/---}' for figures
% '\subfile{\modulefiles/---}' for register files

% >>> M1: This module is in global project? <<<
% 'YES' - uncomment; 'NO' - comment out
\globaltrue

\ifglobal
    \documentclass[main\_\_EN.tex]{subfiles}
   
\else
    \documentclass[10pt]{article}

    % >>> M2: In ./00-shared/preamble-trm-module, also find
    %         settings for visibility of todo notes, watermark <<<
    \usepackage{preamble-trm-module}

    % >>> M3: Temporary labels in English,
    %         you can add your own labels there <<<
    \myexternaldocument{./00-shared/temp-labels-en}

\fi %\ifglobal


\setcounter{secnumdepth}{4}

% >>> M5: If this module needs any updates in
% './00-shared/preamble-trm-module.sty'
% (include additional package, etc.), contact your module PM
% or create the file './00-shared/changelog.tex' make the
% necessary updates yourself and document them in this file <<<

% Make text left-aligned, comment out for full justification
\raggedright


\begin{document}


% Sets language into which pre-defined bits of text are translated
\selectlanguage{english}

% Do not delete this block
\ifglobal
% Add the following front matter if
% this module is outside of global project
\else
    \InputIfFileExists{readme}{}{}
    {\let\clearpage\relax \listoftodos}
    {\let\clearpage\relax \listofdonetodos}
    \tableofcontents
    %\listoftables
    %\listoffigures
\fi


%%%%%%%%%%%%%%%%%%%%%%%%%
%%% TRM Module Begins %%%
%%%%%%%%%%%%%%%%%%%%%%%%%

\hypertarget{mipidsi}{}
\chapter{MIPI DSI}
\label{mod:mipidsi}



%%%%%%%%%%%%%%%%%%%%
%%% Introduction %%%
%%%%%%%%%%%%%%%%%%%%

\section{Introduction}
MIPI DSI stands for Mobile Industry Processor Interface Display Serial Interface. It is a standard interface protocol for connecting displays to mobile devices, such as smartphones, tablets, and other embedded systems. MIPI DSI in \chipseries{} consists of three parts: a DSI Bridge, an MIPI DSI Host Controller, and an MIPI Bidirectional D-PHY. It provides one clock lane and two data lanes, and supports a maximum bit rate of 1.5 Gbps per lane for communication with a DSI-compliant display. DSI (Display Serial Interface) is part of a group of communication protocols defined by the MIPI Alliance.

The DSI Bridge provides an interface between system memory and the MIPI DSI Host Controller. It receives image data through VDMA from system memory, and outputs image data through the Display Pixel Interface (DPI) to the MIPI DSI Host Controller.

The MIPI DSI Host Controller implements the protocol functions defined in the MIPI DSI Specification. It receives image data from the DSI Bridge and outputs image data through the PHY-Protocol Interface (PPI) to the MIPI Bidirectional D-PHY.

The MIPI Bidirectional D-PHY implements the physical layer of bidirectional universal lanes for the MIPI D-PHY interface, which provides one clock lane and two data lanes. The PHY is only used as master.


%%%%%%%%%%%%%%%%%%%%
%%% Terminology %%%
%%%%%%%%%%%%%%%%%%%%



%%%%%%%%%%%%%%%%%%%%
%%% Feature List %%%
%%%%%%%%%%%%%%%%%%%%

\section{Feature List}
\label{sec:mipidsi-features}

MIPI DSI in \chipname{} supports the following features:
\begin{itemize}
    \item Configuration of DSI bridge or VDMA as flow controller when reading image data from memory
    \item Configurable DPI video timing in DSI Bridge
    \item Various input image formats: RGB888, RGB666, RGB565, YUV422
    \item Various DPI image output formats: RGB888, RGB666, RGB565
    \item Maximum DPI frequency: 240 MHz
    \item MIPI DSI Host compliant with MIPI DSI specification
    \item PHY compliant with MIPI D-PHY interface specification, revision 1.1
    \item The PHY includes a PLL
    \item Continuous and non-continuous clock-lane modes
    \item Escape Mode through data lane 0
    \item Lane operation ranging from 80 Mbps to 1.5 Gbps in HS TX
    \item Burst and non-burst Video Modes
    \item Command Mode through APB interface to access display registers
    \item Supports commands defined in MIPI Alliance specification for Display Command Set (DCS)
    \item Video pattern generator
    \item Supports Ultra Low-Power mode
\end{itemize}


%%%%%%%%%%%%%%%%%%%%%%%%%%%%%%
%%% Architectural Overview %%%
%%%%%%%%%%%%%%%%%%%%%%%%%%%%%%

\section{Architectural Overview}
\label{sec:mipidsi-arch-overview}
Figure \ref{figure:mipidsi_arch} shows the architecture of MIPI DSI.

\begin{figure}[H]
    \centering
    \includegraphics[width=1.0\textwidth]{\modulefiles/figures/mipidsi_arch.pdf}
    \caption{MIPI DSI architecture}
    \label{figure:mipidsi_arch}
\end{figure}


\begin{itemize}
    \item \textbf{DSI Bridge}: The DSI Bridge receives image data from system memory through VDMA, converts it to the DPI data format, and then outputs it to the MIPI DSI Host.
    \item \textbf{MIPI DSI Host}: The MIPI DSI Host receives image data from two sources: the DSI Bridge, or the video pattern generator. Then the received image data is packaged by the Packet Handler according to the MIPI DSI Specification and sent to the PHY Interface Controller. Finally, the image data is converted to the PPI data format and output to the MIPI DSI Bidirectional D-PHY. The MIPI DSI Host includes:

    \begin{itemize}
        \item \textbf{DPI Interface}: Receives pixel data from the DSI Bridge and buffers it to a FIFO.
        \item \textbf{Register Bank}: Provides access to configuration and control registers.
        \item \textbf{APB to Generic Interface}: Converts APB operations into Generic Commands.
        \item \textbf{Video Pattern Generator}: Generates video test patterns.
        \item \textbf{Packet Handler}: Schedules activities inside the link, including packet building, generating ECC CRC codes, packet reception, and so on.
        \item \textbf{PHY Interface Controller}: Manages the PHY PPI interfaces, including enabling low-power (LP) transmission/reception, high-speed (HS) transmission, and so on.
        \item \textbf{Error Management}: Monitors and notifies error conditions on the DSI link.
    \end{itemize}
    \item \textbf{MIPI DSI D-PHY}: The MIPI DSI Bidirectional D-PHY implements the physical layer with one clock lane and two data lanes, connecting to the external display driver. MIPI DSI Bidirectional D-PHY also includes a PLL to generate the HS TX clock.
\end{itemize}


%%%%%%%%%%%%%%%%%%%%%%%%%%%%%%
%%% Functional Description %%%
%%%%%%%%%%%%%%%%%%%%%%%%%%%%%%

\section{Functional Description}
\label{sec:mipidsi-func-descr}
This section introduces different operating states and modes of MIPI Bi-Direction (BD) D-PHY and some major functions that facilitate the use of MIPI DSI.

\subsection{DSI Bridge}
\label{sec:dsi-bridge}

DSI Bridge receives image data from system memory through VDMA. The data transfer flow controller can be configured to VDMA or DSI Bridge through \hyperref[fielddesc:DSIBRGDSIDMAFLOWCONTROLLER]{DSI\_BRG\_DSI\_DMA\_FLOW\_CONTROLLER}. It is recommended to use VDMA as the flow controller for better performance.





When receiving image data with VDMA as the flow controller, the image can be split into multiple block transfers. To enable this functionality, set \hyperref[fielddesc:DSIBRGDMAMULTIBLKEN]{DSI\_BRG\_DMA\_MULTIBLK\_EN} to 1, and configure \hyperref[fielddesc:DSIBRGDMAFLOWMULTIBLKNUM]{DSI\_BRG\_DMA\_FLOW\_MULTIBLK\_NUM} with the desired block number. For a detailed explanation of block transfers, refer to \hyperref[sec:vdma-transfer-hierarchy]{Transfer Hierarchy} in Chapter \ref{mod:vdma} \textit{\nameref{mod:vdma}}.

DSI Bridge supports input of RGB888, RGB666, RGB565, and YUV422 data types. For YUV422, DSI Bridge supports 8 different pixel arrangements. Users can select the input data type as RGB or YUV by configuring the \hyperref[fielddesc:DSIBRGDATAINTYPE]{DSI\_BRG\_DATA\_IN\_TYPE} register. If RGB is selected, use \hyperref[fielddesc:DSIBRGRAWTYPE]{DSI\_BRG\_RAW\_TYPE} to specify a specific RGB format. If YUV is selected, also use \hyperref[fielddesc:DSIBRGRAWTYPE]{DSI\_BRG\_RAW\_TYPE} to specify the RGB format to which the YUV data will be converted.

All supported input pixel arrangements are listed in table \ref{tab:input-pixel-arrangement}.


\scriptsize
\begin{longtable}[c]{ |p{2cm}|p{1.5cm}|p{1.5cm}|p{1.5cm}|p{1.5cm}|p{1.5cm}|p{1.5cm}|p{1.5cm}|p{1.5cm}| }
    \caption{Input Pixel Arrangement}
    \label{tab:input-pixel-arrangement}\\\hline

    \rowcolor{lightgray}
    \textbf{Color Mode} & \textbf{Addr + 7} & \textbf{Addr + 6} & \textbf{Addr + 5} & \textbf{Addr + 4} & \textbf{Addr + 3} & \textbf{Addr + 2} & \textbf{Addr + 1} & \textbf{Addr + 0} \\\hline
    \endfirsthead


    \multicolumn{9}{c}{\bfseries \tablename\ \thetable{} -- cont'd from previous page}\\\hline
    \rowcolor{lightgray}
    \textbf{Color Mode} & \textbf{Addr + 7} & \textbf{Addr + 6} & \textbf{Addr + 5} & \textbf{Addr + 4} & \textbf{Addr + 3} & \textbf{Addr + 2} & \textbf{Addr + 1} & \textbf{Addr + 0} \\\hline
    \endhead

    \multicolumn{9}{r}{\textbf{Cont'd on next page}} \\
    \endfoot

    \endlastfoot
        \multirow{3}{*}{\textbf{RGB888}} & \textbf{G2}[7:0] & \textbf{B2}[7:0] & \textbf{R1}[7:0] & \textbf{G1}[7:0] & \textbf{B1}[7:0] & \textbf{R0}[7:0] & \textbf{G0}[7:0] & \textbf{B0}[7:0] \\\cline{2-9}
        & \textbf{B5}[7:0] & \textbf{R4}[7:0] & \textbf{G4}[7:0] & \textbf{B4}[7:0] & \textbf{R3}[7:0] & \textbf{G3}[7:0] & \textbf{B3}[7:0] & \textbf{R2}[7:0] \\\cline{2-9}
        & \textbf{R7}[7:0] & \textbf{G7}[7:0] & \textbf{B7}[7:0] & \textbf{R6}[7:0] & \textbf{G6}[7:0] & \textbf{B6}[7:0] & \textbf{R5}[7:0] & \textbf{G5}[7:0] \\\hline
        \textbf{RGB565} & \textbf{R3}[4:0]\textbf{G3}[5:3] & \textbf{G3}[2:0]\textbf{B3}[4:0] & \textbf{R2}[4:0]\textbf{G2}[5:3] & \textbf{G2}[2:0]\textbf{B2}[4:0]  & \textbf{R1}[4:0]\textbf{G1}[5:3] & \textbf{G1}[2:0]\textbf{B1}[4:0] & \textbf{R0}[4:0]\textbf{G0}[5:3] & \textbf{G0}[2:0]\textbf{B0}[4:0] \\\hline
        \makecell[l]{\textbf{YUV422}\\ \hyperref[fielddesc:DSIBRGYUV422FORMAT]{\makecell[l]{DSI\_BRG\_YUV422\\\_FORMAT = 0}}\\ \hyperref[fielddesc:DSIBRGYUVPIXENDIAN]{\makecell[l]{DSI\_BRG\_YUV\\\_PIX\_ENDIAN = 0}} } & \textbf{Y3}[7:0] & \textbf{U3}[7:0] & \textbf{Y2}[7:0] & \textbf{V2}[7:0] & \textbf{Y1}[7:0] & \textbf{U1}[7:0] & \textbf{Y0}[7:0] & \textbf{V0}[7:0] \\\hline
        \makecell[l]{\textbf{YUV422}\\ \hyperref[fielddesc:DSIBRGYUV422FORMAT]{\makecell[l]{DSI\_BRG\_YUV422\\\_FORMAT = 0}}\\ \hyperref[fielddesc:DSIBRGYUVPIXENDIAN]{\makecell[l]{DSI\_BRG\_YUV\\\_PIX\_ENDIAN = 1}}  } & \textbf{Y0}[7:0] & \textbf{U0}[7:0] & \textbf{Y1}[7:0] & \textbf{V1}[7:0] & \textbf{Y2}[7:0] & \textbf{U2}[7:0] & \textbf{Y3}[7:0] & \textbf{V3}[7:0] \\\hline
        \makecell[l]{\textbf{YUV422}\\ \hyperref[fielddesc:DSIBRGYUV422FORMAT]{\makecell[l]{DSI\_BRG\_YUV422\\\_FORMAT = 1}} \\ \hyperref[fielddesc:DSIBRGYUVPIXENDIAN]{\makecell[l]{DSI\_BRG\_YUV\\\_PIX\_ENDIAN = 0}}  } & \textbf{Y3}[7:0] & \textbf{V3}[7:0] & \textbf{Y2}[7:0] & \textbf{U2}[7:0] & \textbf{Y1}[7:0] & \textbf{V1}[7:0] & \textbf{Y0}[7:0] & \textbf{U0}[7:0] \\\hline
        \makecell[l]{\textbf{YUV422}\\ \hyperref[fielddesc:DSIBRGYUV422FORMAT]{\makecell[l]{DSI\_BRG\_YUV422\\\_FORMAT = 1}} \\ \hyperref[fielddesc:DSIBRGYUVPIXENDIAN]{\makecell[l]{DSI\_BRG\_YUV\\\_PIX\_ENDIAN = 1}}  } & \textbf{Y0}[7:0] & \textbf{V0}[7:0] & \textbf{Y1}[7:0] & \textbf{U1}[7:0] & \textbf{Y2}[7:0] & \textbf{V2}[7:0] & \textbf{Y3}[7:0] & \textbf{U3}[7:0] \\\hline
        \makecell[l]{\textbf{YUV422}\\ \hyperref[fielddesc:DSIBRGYUV422FORMAT]{\makecell[l]{DSI\_BRG\_YUV422\\\_FORMAT = 2}} \\ \hyperref[fielddesc:DSIBRGYUVPIXENDIAN]{\makecell[l]{DSI\_BRG\_YUV\\\_PIX\_ENDIAN = 0}}  } & \textbf{U3}[7:0] & \textbf{Y3}[7:0] & \textbf{V2}[7:0] & \textbf{Y2}[7:0] & \textbf{U1}[7:0] & \textbf{Y1}[7:0] & \textbf{V0}[7:0] & \textbf{Y0}[7:0] \\\hline
        \makecell[l]{\textbf{YUV422}\\ \hyperref[fielddesc:DSIBRGYUV422FORMAT]{\makecell[l]{DSI\_BRG\_YUV422\\\_FORMAT = 2}} \\ \hyperref[fielddesc:DSIBRGYUVPIXENDIAN]{\makecell[l]{DSI\_BRG\_YUV\\\_PIX\_ENDIAN = 1}}  } & \textbf{U0}[7:0] & \textbf{Y0}[7:0] & \textbf{V1}[7:0] & \textbf{Y1}[7:0] & \textbf{U2}[7:0] & \textbf{Y2}[7:0] & \textbf{V3}[7:0] & \textbf{Y3}[7:0] \\\hline
        \makecell[l]{\textbf{YUV422}\\ \hyperref[fielddesc:DSIBRGYUV422FORMAT]{\makecell[l]{DSI\_BRG\_YUV422\\\_FORMAT = 3}} \\ \hyperref[fielddesc:DSIBRGYUVPIXENDIAN]{\makecell[l]{DSI\_BRG\_YUV\\\_PIX\_ENDIAN = 0}} } & \textbf{V3}[7:0] & \textbf{Y3}[7:0] & \textbf{U2}[7:0] & \textbf{Y2}[7:0] & \textbf{V1}[7:0] & \textbf{Y1}[7:0] & \textbf{U0}[7:0] & \textbf{Y0}[7:0] \\\hline
        \makecell[l]{\textbf{YUV422}\\ \hyperref[fielddesc:DSIBRGYUV422FORMAT]{\makecell[l]{DSI\_BRG\_YUV422\\\_FORMAT = 3}} \\ \hyperref[fielddesc:DSIBRGYUVPIXENDIAN]{\makecell[l]{DSI\_BRG\_YUV\\\_PIX\_ENDIAN = 1}}  } & \textbf{V0}[7:0] & \textbf{Y0}[7:0] & \textbf{U1}[7:0] & \textbf{Y1}[7:0] & \textbf{V2}[7:0] & \textbf{Y2}[7:0] & \textbf{U3}[7:0] & \textbf{Y3}[7:0] \\\hline
        

\end{longtable}
\normalsize

DSI Bridge outputs data through the DPI interface, and supports several RGB output bit arrangements specified by \hyperref[fielddesc:DSIBRGDPICONFIG]{DSI\_BRG\_DPI\_CONFIG}. All supported bit arrangements are listed in table \ref{tab:dpi-bit-arrangement}.

\scriptsize
\begin{longtable}[c]{ |p{2.5cm}|p{1cm}|p{1cm}|p{1cm}|p{1cm}|p{1cm}|p{1cm}| }
    \caption{DPI bit arrangement}
    \label{tab:dpi-bit-arrangement}\\\hline

    \rowcolor{lightgray}
    \textbf{} & \multicolumn{3}{|c|}{\textbf{RGB565}} & \multicolumn{2}{|c|}{\textbf{RGB666}} & \textbf{RGB888}  \\\hline
    \rowcolor{lightgray}
    \textbf{\hyperref[fielddesc:DSIBRGDPICONFIG]{DSI\_BRG\_DPI\_CONFIG}} & \textbf{0} & \textbf{1} & \textbf{2} & \textbf{0} & \textbf{1} & \textbf{0}  \\\hline
    \endfirsthead


    \multicolumn{7}{c}{\bfseries \tablename\ \thetable{} -- cont'd from previous page}\\\hline
    \rowcolor{lightgray}
    \textbf{} & \multicolumn{3}{|c|}{\textbf{RGB565}} & \multicolumn{2}{|c|}{\textbf{RGB666}} & \textbf{RGB888}  \\\hline
    \rowcolor{lightgray}
    \textbf{\hyperref[fielddesc:DSIBRGDPICONFIG]{DSI\_BRG\_DPI\_CONFIG}} & \textbf{0} & \textbf{1} & \textbf{2} & \textbf{0} & \textbf{1} & \textbf{0} \\\hline
    \endhead

    \multicolumn{7}{r}{\textbf{Cont'd on next page}} \\
    \endfoot

    \endlastfoot
        
        bit-23 &    &    &    &    &    & R7 \\\hline
        bit-22 &    &    &    &    &    & R6 \\\hline
        bit-21 &    &    & R4 &    & R5 & R5 \\\hline
        bit-20 &    & R4 & R3 &    & R4 & R4 \\\hline
        bit-19 &    & R3 & R2 &    & R3 & R3 \\\hline
        bit-18 &    & R2 & R1 &    & R2 & R2 \\\hline
        bit-17 &    & R1 & R0 & R5 & R1 & R1 \\\hline
        bit-16 &    & R0 &    & R4 & R0 & R0 \\\hline
        bit-15 & R4 &    &    & R3 &    & G7 \\\hline
        bit-14 & R3 &    &    & R2 &    & G6 \\\hline
        bit-13 & R2 & G5 & G5 & R1 & G5 & G5 \\\hline
        bit-12 & R1 & G4 & G4 & R0 & G4 & G4 \\\hline
        bit-11 & R0 & G3 & G3 & G5 & G3 & G3 \\\hline
        bit-10 & G5 & G2 & G2 & G4 & G2 & G2 \\\hline
        bit-9  & G4 & G1 & G1 & G3 & G1 & G1 \\\hline
        bit-8  & G3 & G0 & G0 & G2 & G0 & G0 \\\hline
        bit-7  & G2 &    &    & G1 &    & B7 \\\hline
        bit-6  & G1 &    &    & G0 &    & B6 \\\hline
        bit-5  & G0 &    & B4 & B5 & B5 & B5 \\\hline
        bit-4  & B4 & B4 & B3 & B4 & B4 & B4 \\\hline
        bit-3  & B3 & B3 & B2 & B3 & B3 & B3 \\\hline
        bit-2  & B2 & B2 & B1 & B2 & B2 & B2 \\\hline
        bit-1  & B1 & B1 & B0 & B1 & B1 & B1 \\\hline
        bit-0  & B0 & B0 &    & B0 & B0 & B0 \\\hline
        

\end{longtable}
\normalsize

\subsection{DSI Host}

\subsubsection{Video Transmission}
DSI Host uses DPI as input interface. It is used to transmit a real-time pixel stream from HOST to peripherals in Video Mode.

There are two kinds of video transmission modes:

\begin{itemize}
    \item \textbf{Burst Mode}: In Burst Mode, one pixel line is buffered in FIFO and transmitted in one single packet without any interruptions. It allows the pixel line to be sent in a short time and then return to lane LP state, which can reduce power consumption and enable more time for other data transmission.
    \item \textbf{Non-Burst Mode}: In Non-Burst Mode, one pixel line is divided into several packets. It is used to match pixel bandwidth with DSI link bandwidth, and reduce the memory requirement for both host and peripheral. There are two kinds of Non-Burst mode:
    \begin{itemize}
        \item Non-Burst mode with sync pulse: allows the peripheral to accurately reconstruct original video timing, including sync pulse widths.
        \item Non-Burst mode with sync event: accurate reconstruction of sync pulse is not required.
    \end{itemize}
\end{itemize}

If the host and peripheral meet the following requirements, it is recommended to select Burst Mode. Otherwise, the pixel data may be lost, causing malfunction of the display peripheral.
\begin{itemize}
    \item The display peripheral supports receiving a full pixel line in one single block and the reception buffer will not overflow.
    \item The DSI output bandwidth is much higher than the DPI input bandwidth, which allows the lane to go to the LP state after each line.
\end{itemize}

If any of the requirements cannot be met, Non-Burst Mode should be selected. In Non-Burst Mode, the DSI host should be configured to match DSI output ratio and DPI input ratio. The following equation shows how the transmission parameters should be configured (The left side represents the DSI output ratio, and the right side represents the DPI input ratio):
\begin{itemize}
    \item If null packets are enabled (\hyperref[fielddesc:DSIHOSTVIDNULLSIZE]{DSI\_HOST\_VID\_NULL\_SIZE} != 0):
    $ lanebyteclk\_period \times \hyperref[fielddesc:DSIHOSTVIDNUMCHUNKS]{DSI\_HOST\_VID\_NUM\_CHUNKS} \times  (\hyperref[fielddesc:DSIHOSTVIDPKTSIZE]{DSI\_HOST\_VID\_PKT\_SIZE} \times  bytes\_per\_pixel + 12 + \hyperref[fielddesc:DSIHOSTVIDNULLSIZE]{DSI\_HOST\_VID\_NULL\_SIZE})/number\_of\_lanes = dpiclk\_period \times  pixels\_per\_line $
    \item If null packets are disabled (\hyperref[fielddesc:DSIHOSTVIDNULLSIZE]{DSI\_HOST\_VID\_NULL\_SIZE} = 0):
    $ lanebyteclk\_period \times  \hyperref[fielddesc:DSIHOSTVIDNUMCHUNKS]{DSI\_HOST\_VID\_NUM\_CHUNKS} \times  (\hyperref[fielddesc:DSIHOSTVIDPKTSIZE]{DSI\_HOST\_VID\_PKT\_SIZE} \times  bytes\_per\_pixel + 6)/number\_of\_lanes = dpiclk\_period \times  pixels\_per\_line $
\end{itemize}

Refer to \ref{sec:clocks-and-display-resolution} for the definitions of lanebyteclk and dpiclk.

\subsubsection{Command Transmission}
\label{sec:command-transmission}
DSI Host allows the transmission of generic commands described in DSI specification as manufacturer-specific in Command Mode through APB Slave Interface.

The command packets are built using the APB registers:
\begin{itemize}
    \item \textbf{\hyperref[regdesc:DSIHOSTGENPLDDATAREG]{DSI\_HOST\_GEN\_PLD\_DATA\_REG}}: Writing to this register sends the data as the payload of a command packet. Reading this register returns the payload of read-back packets.
    \item \textbf{\hyperref[regdesc:DSIHOSTGENHDRREG]{DSI\_HOST\_GEN\_HDR\_REG}}: Contains the command packet header type and header data. Writing to this register triggers a command transmission.
\end{itemize}

The valid packets are as follows:
\begin{itemize}
    \item Generic Write Short Packet 0 Parameters
    \item Generic Write Short Packet 1 Parameters
    \item Generic Write Short Packet 2 Parameters
    \item Generic Read Short Packet 0 Parameters
    \item Generic Read Short Packet 1 Parameters
    \item Generic Read Short Packet 2 Parameters
    \item Maximum Read Packet Configuration
    \item Generic Long Write Packet
\end{itemize}

Command packets can be transmitted in lane HS or LP states, while in video mode Blanking or Low-Power periods (BLLP). BLLP corresponds to the shaded areas in figure \ref{figure:command-transmission-periods}.

\begin{figure}[H]
    \centering
    \includegraphics[width=0.8\textwidth]{\modulefiles/figures/command-transmission-periods.pdf}
    \caption{Command Transmission Periods}
    \label{figure:command-transmission-periods}
\end{figure}

\paragraph{Command Transmission in LP mode}

To enable sending commands in LP state during video transmission, set \hyperref[fielddesc:DSIHOSTLPCMDEN]{DSI\_HOST\_LP\_CMD\_EN} to 1. In this case, even if bits 8 to 13 of \hyperref[regdesc:DSIHOSTVIDMODECFGREG]{DSI\_HOST\_VID\_MODE\_CFG\_REG} are set to 0, commands can still be sent during these periods.

\begin{itemize}
    \item \hyperref[fielddesc:DSIHOSTOUTVACTLPCMDTIME]{DSI\_HOST\_OUTVACT\_LPCMD\_TIME} configures the time available (in bytes) to transmit commands in LP mode (based on escape clock) during VSA, VBP, and VFP. The outvact\_lpcmd\_time for Non-Burst with Sync Pulse is shown in figure \ref{figure:outvact-lpcmd-time-for-non-burst-with-sync-pulse}, and the outvact\_lpcmd\_time for Burst or Non-Burst with Sync Event is shown in figure \ref{figure:outvact-lpcmd-time-for-burst-or-non-burst-with-sync-event}.
    \begin{figure}[H]
        \centering
        \includegraphics[width=1.0\textwidth]{\modulefiles/figures/outvact-lpcmd-time-for-non-burst-with-sync-pulse.pdf}
        \caption{outvact\_lpcmd\_time for Non-Burst with Sync Pulse}
        \label{figure:outvact-lpcmd-time-for-non-burst-with-sync-pulse}
    \end{figure}
    \begin{figure}[H]
        \centering
        \includegraphics[width=1.0\textwidth]{\modulefiles/figures/outvact-lpcmd-time-for-burst-or-non-burst-with-sync-event.pdf}
        \caption{outvact\_lpcmd\_time for Burst or Non-Burst with Sync Event}
        \label{figure:outvact-lpcmd-time-for-burst-or-non-burst-with-sync-event}
    \end{figure}
    The outvact\_lpcmd\_time is calculated by:\\
    $outvact\_lpcmd\_time = (t_L - (t_{H1} + t_{HS->LP} + t_{LP->HS} + t_{LPDT} + 2t_{ESCCLK})) / (2 \times 8 \times t_{ESCCLK})$
    
    \item \hyperref[fielddesc:DSIHOSTINVACTLPCMDTIME]{DSI\_HOST\_INVACT\_LPCMD\_TIME} configures the time available (in bytes) to transmit commands in LP mode (based on escape clock) during HFP. The invact\_lpcmd\_time for Non-Burst with Sync Pulse is shown in figure \ref{figure:invact-lpcmd-time-for-non-burst-with-sync-pulse}; and the invact\_lpcmd\_time for Non-Burst with Sync Event is shown in figure \ref{figure:invact-lpcmd-time-for-non-burst-with-sync-event}.
    \begin{figure}[H]
        \centering
        \includegraphics[width=1.0\textwidth]{\modulefiles/figures/invact-lpcmd-time-for-non-burst-with-sync-pulse.pdf}
        \caption{invact\_lpcmd\_time for Non-Burst with Sync Pulse}
        \label{figure:invact-lpcmd-time-for-non-burst-with-sync-pulse}
    \end{figure}
    \begin{figure}[H]
        \centering
        \includegraphics[width=1.0\textwidth]{\modulefiles/figures/invact-lpcmd-time-for-non-burst-with-sync-event.pdf}
        \caption{invact\_lpcmd\_time for Burst or Non-Burst with Sync Event}
        \label{figure:invact-lpcmd-time-for-non-burst-with-sync-event}
    \end{figure}
    The invact\_lpcmd\_time is calculated by:\\
    $invact\_lpcmd\_time = (t_L - (t_{HSA} + t_{HBP} + t_{HACT} + t_{HS->LP} + t_{LP->HS} + t_{LPDT} + 2t_{ESCCLK})) / (2 \times 8 \times t_{ESCCLK})$
\end{itemize}


\paragraph{Command Transmission in HS Mode}
If \hyperref[fielddesc:DSIHOSTLPCMDEN]{DSI\_HOST\_LP\_CMD\_EN} is set to 0, the commands are sent in HS mode during video transmission. In this case, no programming is required.

\subsubsection{Video Pattern Generator}
\label{sec:video-pattern-generator}
The Video Pattern Generator can output horizontal or vertical color bar or BER testing pattern. For the pattern resolution, refer to \hyperref[sec:clocks-and-display-resolution]{Clocks and Display Resolution}. The following register fields are used to configure the Video Pattern Generator:

\begin{itemize}
    \item \hyperref[fielddesc:DSIHOSTVPGORIENTATION]{DSI\_HOST\_VPG\_ORIENTATION}: Indicates the orientation of the color bar
    \item \hyperref[fielddesc:DSIHOSTVPGMODE]{DSI\_HOST\_VPG\_MODE}: Selects the pattern
    \item \hyperref[fielddesc:DSIHOSTVPGEN]{DSI\_HOST\_VPG\_EN}: Enables the video pattern generator
\end{itemize}

\paragraph{Color Bar Pattern}
There are eight colors for the color bars: white, yellow, cyan, green, magenta, red, blue, and black. Each color width can be calculated by dividing the pattern size by eight. If the size is not divisible by eight, the last color black contains the remaining size. Figure \ref{figure:vertical-color-bar} and \ref{figure:horizontal-color-bar} show the vertical and horizontal color bars.

\begin{figure}[H]
    \centering
    \includegraphics[width=0.6\textwidth]{\modulefiles/figures/vertical-color-bar.pdf}
    \caption{Vertical Color Bar}
    \label{figure:vertical-color-bar}
\end{figure}

\begin{figure}[H]
    \centering
    \includegraphics[width=0.6\textwidth]{\modulefiles/figures/horizontal-color-bar.pdf}
    \caption{Horizontal Color Bar}
    \label{figure:horizontal-color-bar}
\end{figure}

Table \ref{tab:color-bar-coding} shows the color bar pattern RGB coding.
\begin{longtable}[c]{ |p{0.5cm}|p{1.5cm}|p{1.5cm}|p{1.5cm}|p{1.5cm}|p{1.5cm}|p{1.5cm}|p{1.5cm}|p{1.5cm}| }
    \caption{Color Bar Coding}
    \label{tab:color-bar-coding}\\\hline

    \rowcolor{lightgray}
     & \textbf{White} & \textbf{Yellow} & \textbf{Cyan} & \textbf{Green} & \textbf{Magenta} & \textbf{Red} & \textbf{Blue} & \textbf{Black}  \\\hline
    \endfirsthead


    \multicolumn{9}{c}{\bfseries \tablename\ \thetable{} -- cont'd from previous page}\\\hline
    \rowcolor{lightgray}
    & \textbf{White} & \textbf{Yellow} & \textbf{Cyan} & \textbf{Green} & \textbf{Magenta} & \textbf{Red} & \textbf{Blue} & \textbf{Black}  \\\hline
    \endhead

    \multicolumn{9}{r}{\textbf{Cont'd on next page}} \\
    \endfoot

    \endlastfoot
        
        \textbf{R} & High & High & Low  & Low  & High & High & Low  & Low \\\hline
        \textbf{G} & High & High & High & High & Low  & Low  & Low  & Low \\\hline
        \textbf{B} & High & Low  & High & Low  & High & Low  & High & Low \\\hline

\end{longtable}

\paragraph{BER Testing Pattern}
The BER testing pattern tests whether the device RX PHY receives data correctly. The following data patterns are generated:
\begin{itemize}
    \item X bytes of 0xAA (high-frequency pattern)
    \item X bytes of 0x33 (mid-frequency pattern)
    \item X bytes of 0xF0 (low-frequency pattern)
    \item X bytes of 0x7F (lone 0 pattern)
    \item X bytes of 0x55 (high-frequency pattern)
    \item X bytes of 0xCC (mid-frequency pattern)
    \item X bytes of 0x0F (low-frequency pattern)
    \item Y bytes of 0x80 (lone 1 pattern)
\end{itemize}

X and Y depend on pattern resolution, and Y equals X in most cases. With RGB888 color coding and horizontal resolution in multiples of eight, Figure~\ref{figure:rgb888-ber-testing-pattern} shows the pattern that appears on the DSI display:

\begin{figure}[H]
    \centering
    \includegraphics[width=0.6\textwidth]{\modulefiles/figures/rgb888-ber-testing-pattern.pdf}
    \caption{RGB888 BER Testing Pattern}
    \label{figure:rgb888-ber-testing-pattern}
\end{figure}

\subsubsection{Error Reporting}
The \hyperref[regdesc:DSIHOSTINTST0REG]{DSI\_HOST\_INT\_ST\textcolor{red}{0/1}\_REG} registers indicate the DSI Host error reports. These registers can also trigger an interrupt to notify the system. To trigger the error interrupt, the corresponding bit in mask registers \hyperref[regdesc:DSIHOSTINTMSK0REG]{DSI\_HOST\_INT\_MSK\textcolor{red}{0/1}\_REG} should be cleared. The \hyperref[regdesc:DSIHOSTINTST0REG]{DSI\_HOST\_INT\_ST\textcolor{red}{0/1}\_REG} will be cleared after a read operation.





\subsection{D-PHY}
This section describes the D-PHY operating modes and PLL configurations.

\subsubsection{D-PHY Operating Modes}
Figure \ref{figure:phy-operating-modes} shows the PHY initialization and active operating modes.
\begin{figure}[H]
    \centering
    \includegraphics[width=1.0\textwidth]{\modulefiles/figures/phy-operating-modes.pdf}
    \caption{PHY operating modes}
    \label{figure:phy-operating-modes}
\end{figure}

\paragraph{No-Power State}
No-Power state is characterized by the absence of any supply voltage applied to the PHY.

For \chipname{}'s MIPI DSI, MIPI BD D-PHY is in No-Power state when the power supply VDD\_HP\_\regindex{n} (n: 0-3) and VDD\_MIPI\_DPHY are off.

\paragraph{Shutdown State}
In Shutdown state, all analog blocks are disabled, all digital logic is reset, and DP/DN and CP/CN lane are high impedance (Hi-Z). To enable shutdown mode of the PHY, clear \hyperref[fielddesc:DSIHOSTPHYRSTZ]{DSI\_HOST\_PHY\_RSTZ} and \hyperref[fielddesc:DSIHOSTPHYSHUTDOWNZ]{DSI\_HOST\_PHY\_SHUTDOWNZ} to 0.

Before leaving the Shutdown state, some parameters should be configured, including the number of lanes to be used, the lane operation frequency range, and the PLL configurations.

The MIPI BD D-PHY consists of two data lanes. Configure \hyperref[fielddesc:DSIHOSTNLANES]{DSI\_HOST\_N\_LANES} to select either one or both data lanes for use.

By default, the PHY is set to operate within a lower range of 80-110 Mbps. To accommodate higher bit rates, you can adjust the HS frequency ranges (the hsfreqrange field) in test code \hyperref[sec:test-code-0x44]{test code 0x44} with the proper code. It is recommended to perform this adjustment in Shutdown state if the PHY works always at the same bit rate, as the control interface is independent of the rest of the MIPI BD D-PHY. The supported ranges are shown in table \ref{tab:mipidsi-hsfreq}:

\begin{longtable}{| p{5cm} | R{5cm} |}
    \caption{Frequency Ranges Supported by MIPI BD D-PHY} \label{tab:mipidsi-hsfreq} \\
    \hline
    
    \rowcolor{lightgray}
    \textbf{Frequency Range (Mbps)} & \textbf{hsfreqrange[5:0]} \\ \hline
    \endfirsthead

    \multicolumn{2}{c}{\bfseries \tablename\ \thetable{} -- continued from previous page}\\\hline
    \rowcolor{lightgray}
    \textbf{Frequency Range (Mbps)} & \textbf{hsfreqrange[5:0]} \\ \hline
    \endhead

    \multicolumn{2}{r}{\textbf{Continued on next page}} \\
    \endfoot

    \endlastfoot
    
    80-89 & 000000  \\ \hline
    90-99 & 010000  \\ \hline
    100-109 & 100000  \\ \hline
    110-129 & 000001  \\ \hline
    130-139 & 010001  \\ \hline
    140-149 & 100001  \\ \hline
    150-169 & 000010  \\ \hline
    170-179 & 010010  \\ \hline
    180-199 & 100010  \\ \hline
    200-219 & 000011  \\ \hline
    220-239 & 010011  \\ \hline
    240-249 & 100011  \\ \hline
    250-269 & 000100  \\ \hline
    270-299 & 010100  \\ \hline
    300-329 & 000101  \\ \hline
    330-359 & 010101  \\ \hline
    360-399 & 100101  \\ \hline
    400-449 & 000110  \\ \hline
    450-499 & 010110  \\ \hline
    500-549 & 000111  \\ \hline
    550-599 & 010111  \\ \hline
    600-649 & 001000  \\ \hline
    650-699 & 011000  \\ \hline
    700-749 & 001001  \\ \hline
    750-799 & 011001  \\ \hline
    800-849 & 101001  \\ \hline
    850-899 & 111001  \\ \hline
    900-949 & 001010  \\ \hline
    950-999 & 011010  \\ \hline
    1000-1049 & 101010  \\ \hline
    1050-1099 & 111010  \\ \hline
    1100-1149 & 001011  \\ \hline
    1150-1199 & 011011  \\ \hline
    1200-1249 & 101011  \\ \hline
    1250-1299 & 111011  \\ \hline
    1300-1349 & 001100  \\ \hline
    1350-1399 & 011100  \\ \hline
    1400-1449 & 101100  \\ \hline
    1450-1500 & 111100  \\ \hline
\end{longtable}

It is also necessary to configure PLL input and output frequency through test code interface. Refer to \hyperref[sec:pll-parameters]{PLL parameters} for detailed description of PLL parameters.

When \hyperref[fielddesc:DSIHOSTPHYRSTZ]{DSI\_HOST\_PHY\_RSTZ} and \hyperref[fielddesc:DSIHOSTPHYSHUTDOWNZ]{DSI\_HOST\_PHY\_SHUTDOWNZ} are set to 1, the MIPI BD D-PHY leaves the Shutdown state and starts an initialization procedure.

\paragraph{PLL-Locking and AFE Initialization}
After \hyperref[fielddesc:DSIHOSTPHYRSTZ]{DSI\_HOST\_PHY\_RSTZ} and \hyperref[fielddesc:DSIHOSTPHYSHUTDOWNZ]{DSI\_HOST\_PHY\_SHUTDOWNZ} are released, the PHY begins an initialization sequence to enable its normal operation. It is recommended to release \hyperref[fielddesc:DSIHOSTPHYSHUTDOWNZ]{DSI\_HOST\_PHY\_SHUTDOWNZ} before releasing \hyperref[fielddesc:DSIHOSTPHYRSTZ]{DSI\_HOST\_PHY\_RSTZ}. Additionally, ensure that DSI\_DPHY\_CFG\_CLK and DSI\_DPHY\_PLL\_REFCLK (refer to Table \ref{tab:resclk-periph-clks-hp2} in Chapter \ref{mod:resclk} \textit{\nameref{mod:resclk}} for details) are available and stable at this point.

After that, the PLL becomes active and the PHY waits for the PLL lock which is indicated by \hyperref[fielddesc:DSIHOSTPHYLOCK]{DSI\_HOST\_PHY\_LOCK}.

Following the PLL lock, AFE Initialization is performed. All initialization steps are done when \hyperref[fielddesc:DSIHOSTPHYSTOPSTATECLKLANE]{DSI\_HOST\_PHY\_STOPSTATECLKLANE} \hyperref[fielddesc:DSIHOSTPHYSTOPSTATE0LANE]{DSI\_HOST\_PHY\_STOPSTATE\textcolor{red}{n}LANE} are asserted and the lanes are in stop state (LP-11).

\paragraph{Control Mode}
After the AFE initialization is completed, the MIPI BD D-PHY remains in Control mode by default while driving the stop state LP-11 on the lines until a request is made. Any request must start from and end in this state. After a request, a lane can enter High-Speed Data Transfer mode, Escape mode, Ultra Low Power state, or turnaround.



\paragraph{High-Speed Data Transfer}
DSI HOST sends a request for High-Speed Data Transfer after receiving data from DSI Bridge. In High-Speed Data Transfer mode, the PHY transmits image data in bursts. The lane operates in high-speed mode only during these bursts.

As shown in Figure \ref{fig:mipidsi-hs}, a burst on lane consists of three parts: the low-power initialization sequence, the high-speed data payload, and the end of transfer sequence.

\begin{figure}[H]
     \centering
     \includegraphics[width=1.0\linewidth]{\modulefiles/figures/mipidsi-hs.pdf}
     \caption{MIPI DSI HS Data Transfer Sequence}
     \label{fig:mipidsi-hs}
 \end{figure}

\paragraph{Escape Mode}
Escape mode allows asynchronous communication over data lanes at low speed. During Escape mode, the lane remains in low-power state. MIPI DSI in \chipname{} supports Low-Power Data Transmission (LPDT) and Ultra Low Power State (ULPS) in Escape mode.

Commands will be sent in LPDT if \hyperref[fielddesc:DSIHOSTLPCMDEN]{DSI\_HOST\_LP\_CMD\_EN} is set to 1.

You can drive the PHY into ULPS through \hyperref[regdesc:DSIHOSTPHYULPSCTRLREG]{DSI\_HOST\_PHY\_ULPS\_CTRL\_REG} when DSI is idle to save power. The lane state can be observed through \hyperref[regdesc:DSIHOSTPHYSTATUSREG]{DSI\_HOST\_PHY\_STATUS\_REG}.

\paragraph{Turnaround}
The PHY can swap transmission direction through turnaround. This is used to receive command response from DSI devices.

\subsubsection{PLL Parameters}
\label{sec:pll-parameters}
The VCO output frequency of the PLL is a function of the frequency of input reference clock DSI\_DPHY\_PLL\_REFCLK, as well as the multiplication ratio $M$, and the division ratio $N$:

\[
f_{VCO} = \frac{M}{N} \times f_{REFCLK}
\]

And the follow limit must be met:

\[
40 MHz \geq \frac{f_{REFCLK}}{N} \geq 5 MHz
\]

Additionally, M must be an even number. The values of $M$ and $N$ can be programmed through \hyperref[sec:test-code-0x17]{test code 0x17} and \hyperref[sec:test-code-0x18]{test code 0x18}. The programmed value is equal to the actual value minus 1, as shown in table \ref{tab:m-n-ratios}. After that, bit-5 and bit-4 of \hyperref[sec:test-code-0x19]{test code 0x19} should be set to 1 to make the M N configuration valid.

\begin{longtable}{|p{3cm}|p{2cm}|p{3cm}|p{2cm}|}
    \caption{M and N ratios} \label{tab:m-n-ratios} \\
    \hline
    \rowcolor{lightgray}
    \textbf{m[8:0] = M-1} & \textbf{M} & \textbf{n[6:0] = N-1} & \textbf{N} \\ \hline
    \endfirsthead

    \multicolumn{4}{c}{\bfseries \tablename\ \thetable{} -- continued from previous page}\\\hline
    \rowcolor{lightgray}
    \textbf{m[8:0] = M-1} & \textbf{M} & \textbf{n[6:0] = N-1} & \textbf{N} \\ \hline
    \endhead

    \multicolumn{4}{r}{\textbf{Continued on next page}} \\
    \endfoot

    \endlastfoot
    
    9'h1 & 2 & 7'h0 & 1  \\ \hline
    9'h3 & 4 & 7'h1 & 2  \\ \hline
    ... & ... & ... & ...  \\ \hline
\end{longtable}

The VCO range of $f_{VCO}$ can also be selected through \hyperref[sec:test-code-0x10]{test code 0x10}. The valid VCO ranges are in table \ref{tab:vco-ranges}.

\begin{longtable}{|p{3cm}|p{5cm}|}
    \caption{VCO ranges} \label{tab:vco-ranges} \\
    \hline
    \rowcolor{lightgray}
    \textbf{vcorange[2:0]} & \textbf{VCO Frequency (MHz)} \\ \hline
    \endfirsthead

    \multicolumn{2}{c}{\bfseries \tablename\ \thetable{} -- continued from previous page}\\\hline
    \rowcolor{lightgray}
    \textbf{vcorange[2:0]} & \textbf{VCO Frequency (MHz)} \\ \hline
    \endhead

    \multicolumn{2}{r}{\textbf{Continued on next page}} \\
    \endfoot

    \endlastfoot
    
    000 & 80-200  \\ \hline
    001 & 200-300  \\ \hline
    010 & 300-500  \\ \hline
    011 & 500-700  \\ \hline
    100 & 700-900  \\ \hline
    101 & 900-1100  \\ \hline
    110 & 1100-1300  \\ \hline
    111 & 1300-1500  \\ \hline
\end{longtable}

The loop bandwidth should be configured through test code interface depending on the output frequency to ensure proper operation of the PLL. The related parameters are listed below, and the valid value are listed in table \ref{tab:pll-cp-and-lpf-control-bits}.

\begin{itemize}
    \item VCO capacitance (vcocap[1:0]), program through \hyperref[sec:test-code-0x10]{test code 0x10}.
    \item CP current (icpctrl[3:0]), program through \hyperref[sec:test-code-0x11]{test code 0x11}.
    \item LPF characteristic (lpfctrl[5:0]), program through \hyperref[sec:test-code-0x12]{test code 0x12}.
\end{itemize}

\begin{longtable}{|p{3cm}|p{3cm}|p{3cm}|p{2cm}|p{2cm}|p{2cm}|}
    \caption{PLL CP and LPF Control Bits} \label{tab:pll-cp-and-lpf-control-bits} \\
    \hline
    \rowcolor{lightgray}
    \textbf{VCO fmin (MHz)} & \textbf{VCO fmax (MHz)} & \textbf{vcorange[2:0]} & \textbf{icpctrl[3:0]} & \textbf{lpfctrl[5:0]} & \textbf{vcocap[1:0]} \\ \hline
    \endfirsthead

    \multicolumn{6}{c}{\bfseries \tablename\ \thetable{} -- continued from previous page}\\\hline
    \rowcolor{lightgray}
    \textbf{VCO fmin (MHz)} & \textbf{VCO fmax (MHz)} & \textbf{vcorange[2:0]} & \textbf{icpctrl[3:0]} & \textbf{lpfctrl[5:0]} & \textbf{vcocap[1:0]} \\ \hline
    \endhead

    \multicolumn{6}{r}{\textbf{Continued on next page}} \\
    \endfoot

    \endlastfoot
    
    80   & 110  & 000 & 0100 & 010000 & 00 \\ \cline{1-2} \cline{4-6}
    110  & 150  &     & 1000 & 001000 & 00 \\ \cline{1-2} \cline{4-6}
    150  & 200  &     & 1100 & 001000 & 00 \\ \hline
    200  & 250  & 001 & 1100 & 001000 & 00 \\ \cline{1-2} \cline{4-6}
    250  & 300  &     & 0010 & 001000 & 00 \\ \hline
    300  & 400  & 010 & 0101 & 001000 & 00 \\ \cline{1-2} \cline{4-6}
    400  & 500  &     & 0101 & 000001 & 00 \\ \hline
    500  & 600  & 011 & 0110 & 010000 & 00 \\ \cline{1-2} \cline{4-6}
    600  & 700  &     & 0110 & 000100 & 00 \\ \hline
    700  & 900  & 100 & 0110 & 000100 & 00 \\ \hline
    900  & 1100 & 101 & 0111 & 010000 & 00 \\ \hline
    1100 & 1300 & 110 & 0111 & 001000 & 00 \\ \hline
    1300 & 1500 & 111 & 0111 & 001000 & 00 \\ \hline
\end{longtable}

\begin{tiplisting}
    vcorange, icpctrl, lpfctrl are automatically set to the default values for the corresponding hsfreqrange, if the corresponding bit is 0.
    \begin{itemize}
        \item \hyperref[sec:test-code-0x10]{test code 0x10, bit7}, clear to 0 to bypass vcorange
        \item \hyperref[sec:test-code-0x12]{test code 0x12, bit7}, clear to 0 to bypass icpctrl
        \item \hyperref[sec:test-code-0x12]{test code 0x12, bit6}, clear to 0 to bypass lpfctrl
    \end{itemize}
\end{tiplisting}

\subsubsection{Test Codes}
The MIPI BD D-PHY provides a set of test codes for configuring MIPI BD D-PHY operation. The configuration can be done through Test Code Interface using \hyperref[regdesc:DSIHOSTPHYTSTCTRL0REG]{DSI\_HOST\_PHY\_TST\_CTRL0\_REG} and \hyperref[regdesc:DSIHOSTPHYTSTCTRL1REG]{DSI\_HOST\_PHY\_TST\_CTRL1\_REG}. For details of the Test Code Interface, refer to section \ref{sec:mipicsi-rx-dphy-program-interface} in Chapter \ref{mod:mipicsi} \textit{\nameref{mod:mipicsi}}.

All supported Test Codes are:

\paragraph{Normal Operation}
\label{sec:normal-operation}
This mode represents the normal operation of the PHY and also the wake up state of the test interface. In this mode, the test interface is inactive.

\begin{itemize}
    \item \textbf{Test Code}: 0x00
    \item \textbf{Test Data}: None
\end{itemize}

\paragraph{PLL Bias Current Selector / Filter Capacitance Control / VCO Control}
\label{sec:test-code-0x10}
Controls the biasing of the PLL and the VCO transfer function.

\begin{itemize}
    \item \textbf{Test Code}: 0x10
    \item \textbf{Test Dout} 
        \begin{longtable}{ |p{3cm}|p{3cm}|p{3cm}|p{3cm}|p{3cm}|}
            \caption{Test Data of 0x10}
            \label{tab:test-data-0x10} \\
            \hline
            \rowcolor{lightgray}
            \textbf{w-1'b0} & \textbf{w-1'b0} & \textbf{w-3'b000} & \textbf{w-2'b00} & \textbf{w-1'b1} \\ \hline
            Bypass VCO range & Reserved & VCO range control & VCO internal capacitance control & Reference bias current selector  \\ \hline
        \end{longtable}
        \begin{itemize}
            \item Bit 7: Bypass VCO range
                \begin{itemize}
                    \item 0 - VCO range is programmed with the default value for the corresponding hsfreqrange
                    \item 1 - VCO range is programmed with bits 5-3
                \end{itemize}
            \item Bit 6: Reserved
            \item Bits 5-3: VCO range control, refer to table \ref{tab:vco-ranges} for details
            \item Bits 2-1: VCO internal capacitance control
                \begin{itemize}
                    \item 00 - Default capacitance
                    \item 01 - Low capacitance (1/4 of the default value)
                    \item 10 - High capacitance (2 times the default value)
                    \item 11 - Not allowed
                \end{itemize}
            \item Bit 0: Reference bias current selector
                \begin{itemize}
                    \item 0 - Internal biasing
                    \item 1 - Clean external current source of 20uA
                \end{itemize}
        \end{itemize}
    \item \textbf{Test Dout}
        \begin{longtable}{ |p{3cm}|p{3cm}|p{3cm}|p{3cm}|p{3cm}|}
            \caption{Test Dout of 0x10}
            \label{tab:test-dout-0x10} \\
            \hline
            \rowcolor{lightgray}
            \textbf{r-1'b0} & \textbf{r-1'b0} & \textbf{r-3'b000} & \textbf{r-2'b00} & \textbf{r-1'b1} \\ \hline
            Bypass VCO range loopback & 1'b0 & VCO range control loopback & VCO internal capacitance control loopback & Reference bias current selector loopback  \\ \hline
        \end{longtable}
\end{itemize}

\paragraph{PLL CP Control/PLL Lock Bypass for Initialization and for ULP}
\label{sec:test-code-0x11}
Controls the Charge-Pump (CP) current and the bypass of the PLL lock for initialization and ULP exit.

\begin{itemize}
    \item \textbf{Test Code}: 0x11
    \item \textbf{Test Dout} 
        \begin{longtable}{ |p{6cm}|p{6cm}|}
            \caption{Test Data of 0x11}
            \label{tab:test-data-0x11} \\
            \hline
            \rowcolor{lightgray}
            \textbf{w-4'b0000} & \textbf{w-4'b0010} \\ \hline
            PLL lock bypass & CP current  \\ \hline
        \end{longtable}
        \begin{itemize}
            \item Bits 7-4: PLL lock bypass
                \begin{itemize}
                    \item xxx1 - Bypass PLL lock during initialization
                    \item x11x - Bypass PLL lock after exiting ULP
                \end{itemize}
            \item Bits 3-0: CP current (uA)
                \begin{itemize}
                    \item 0000 - 1.5
                    \item 0001 - 3
                    \item 0010 - 4.5
                    \item 0011 - 6
                    \item 0100 - 2.5
                    \item 0101 - 5
                    \item 0110 - 7.5
                    \item 0111 - 10
                    \item 1000 - 3
                    \item 1001 - 6
                    \item 1010 - 9
                    \item 1011 - 12
                    \item 1100 - 4
                    \item 1101 - 8
                    \item 1110 - 12
                    \item 1111 - 16
                \end{itemize}
        \end{itemize}
    \item \textbf{Test Dout}
        \begin{longtable}{ |p{6cm}|p{6cm}|}
            \caption{Test Dout of 0x11}
            \label{tab:test-dout-0x11} \\
            \hline
            \rowcolor{lightgray}
            \textbf{r-1'b0} & \textbf{r-1'b0} \\ \hline
            PLL lock bypass loopback & CP current loopback \\ \hline
        \end{longtable}
\end{itemize}

\paragraph{PLL LPF and CP control}
\label{sec:test-code-0x12}
Controls the Charge-Pump (CP) current and the Low-Pass-Filter (LPF) resistor.
\begin{itemize}
    \item \textbf{Test Code}: 0x12
    \item \textbf{Test Dout} 
        \begin{longtable}{ |p{4cm}|p{4cm}|p{4cm}|}
            \caption{Test Data of 0x12}
            \label{tab:test-data-0x12} \\
            \hline
            \rowcolor{lightgray}
            \textbf{w-1'b0} & \textbf{w-1'b0} & \textbf{w-6'b000010} \\ \hline
            Bypass CP default values & Bypass LPF default values & LPF resistor  \\ \hline
        \end{longtable}
        \begin{itemize}
            \item Bit 7: Bypass CP default values
                \begin{itemize}
                    \item 0 - CP is programmed with the default value for the corresponding hsfreqrange
                    \item 1 - CP is programmed with bits 3-0 using test code 0x11
                \end{itemize}
            \item Bit 6: Bypass LPF default values
                \begin{itemize}
                    \item 0 - LPF is programmed with the default value for the corresponding hsfreqrange
                    \item 1 - LPF is programmed with bits 5-0
                \end{itemize}
            \item Bits 5-0: LPF resistor ($k\Omega$)
                \begin{itemize}
                    \item 000000 - 18
                    \item 000001 - 15.6
                    \item 000010 - 15
                    \item 000100 - 14.4
                    \item 001000 - 12.8
                    \item 010000 - 11.4
                    \item 100000 - 10.5
                    \item others - reserved
                \end{itemize}
        \end{itemize}
    \item \textbf{Test Dout}
        \begin{longtable}{ |p{4cm}|p{4cm}|p{4cm}|}
            \caption{Test Dout of 0x12}
            \label{tab:test-dout-0x12} \\
            \hline
            \rowcolor{lightgray}
            \textbf{r-1'b0} & \textbf{r-1'b0} & \textbf{r-6'b000010} \\ \hline
            CP bypass loopback & LPF bypass loopback & LPF resistor loopback \\ \hline
        \end{longtable}
\end{itemize}


\paragraph{PLL Input Divider Ratio}
\label{sec:test-code-0x17}
Controls the PLL's input divider N. Where $N = n + 1$, $N$ is the actual divider ratio and $n$ is the programmed value in the test code.

To enable the programmed value, ensure that bit 4 of \hyperref[sec:test-code-0x19]{test code 0x19} is set to 1.

\begin{itemize}
    \item \textbf{Test Code}: 0x17
    \item \textbf{Test Dout} 
        \begin{longtable}{ |p{6cm}|p{6cm}|}
            \caption{Test Data of 0x17}
            \label{tab:test-data-0x17} \\
            \hline
            \rowcolor{lightgray}
            \textbf{w-1'b0} & \textbf{w-7'b0000001} \\ \hline
            Reserved & Input divider  \\ \hline
        \end{longtable}
    \item \textbf{Test Dout}
        \begin{longtable}{ |p{6cm}|p{6cm}|}
            \caption{Test Dout of 0x17}
            \label{tab:test-dout-0x17} \\
            \hline
            \rowcolor{lightgray}
            \textbf{r-1'b0} & \textbf{r-7'b0000001} \\ \hline
            Reserved & Input divider loopback \\ \hline
        \end{longtable}
\end{itemize}

\paragraph{PLL Loop Divider Ratio}
\label{sec:test-code-0x18}
Controls the PLL's feedback divider M. $M = m+1$, where $M$ is the actual feedback ratio, and $m$ is the programmed value in the test code.

To enable the programmed value, ensure that bit-5 of \hyperref[sec:test-code-0x19]{test code 0x19} is set to 1.

\begin{itemize}
    \item \textbf{Test Code}: 0x18
    \item \textbf{Test Dout} 
        \begin{longtable}{ |p{4cm}|p{4cm}|p{4cm}|}
            \caption{Test Data of 0x18}
            \label{tab:test-data-0x18} \\
            \hline
            \rowcolor{lightgray}
            \textbf{w-1'b0} & \textbf{w-2'b00} & \textbf{w-5'b00101} \\ \hline
            Set 0 & Reserved & Feedback divider [4:0]  \\ \hline
            \rowcolor{lightgray}
            \textbf{w-1'b1} & \textbf{w-3'b00} & \textbf{w-4'b0000} \\ \hline
            Set 1 & Reserved & Feedback divider [8:5]  \\ \hline
        \end{longtable}
        \begin{itemize}
            \item Bit 7: Program selector
                \begin{itemize}
                    \item 0 - Set 0 is programmed
                    \item 1 - Set 1 is programmed
                \end{itemize}
            \item Bits 6-5: Reserved
            \item Bits 4-0: Feedback divider
        \end{itemize}
    \item \textbf{Test Dout}\\
        TESTDIN[7] - Set 0
        \begin{longtable}{ |p{4cm}|p{8cm}|}
            \caption{Test Dout of 0x18}
            \label{tab:test-dout-0x18-set0} \\
            \hline
            \rowcolor{lightgray}
            \textbf{r-3'b000} & \textbf{r-5'b00101} \\ \hline
            3'b000 & Feedback divider [4:0] loopback \\ \hline
        \end{longtable}
        TESTDIN[7] - Set 1
        \begin{longtable}{ |p{4cm}|p{8cm}|}
            \caption{Test Dout of 0x18}
            \label{tab:test-dout-0x18-set1} \\
            \hline
            \rowcolor{lightgray}
            \textbf{r-4'b0000} & \textbf{r-4'b0000} \\ \hline
            4'b0000 & Feedback divider [8:5] loopback \\ \hline
        \end{longtable}
\end{itemize}

\paragraph{PLL Input and Loop Divider Ratios Control}
\label{sec:test-code-0x19}
Controls the PLL input and loop divider ratios.

\begin{itemize}
    \item \textbf{Test Code}: 0x19
    \item \textbf{Test Dout} 
        \begin{longtable}{ |p{2cm}|p{5cm}|p{5cm}|p{2cm}|}
            \caption{Test Data of 0x19}
            \label{tab:test-data-0x19} \\
            \hline
            \rowcolor{lightgray}
            \textbf{w-2'b00} & \textbf{w-1'b0} & \textbf{w-1'b0} & \textbf{w-4'b0000} \\ \hline
            Reserved & Bypass PLL loop divider default values & Bypass PLL input divider default values & Reserved \\ \hline
        \end{longtable}
        \begin{itemize}
            \item Bits 7-6: Reserved
            \item Bit 5: Bypass the PLL loop divider default values
                \begin{itemize}
                    \item 0 - PLL loop divider is programmed with the default values for the corresponding hsfreqrange
                    \item 1 - PLL loop divider is programmed using \hyperref[sec:test-code-0x18]{test code 0x18}
                \end{itemize}
            \item Bit 4: Bypass the PLL input divider default values
                \begin{itemize}
                    \item 0 - PLL input divider is programmed with the default values for the corresponding hsfreqrange
                    \item 1 - PLL input divider is programmed using \hyperref[sec:test-code-0x17]{test code 0x17}
                \end{itemize}
            \item Bits 3-0: Reserved
        \end{itemize}
    \item \textbf{Test Dout}
        \begin{longtable}{ |p{2cm}|p{5cm}|p{5cm}|p{2cm}|}
            \caption{Test Dout of 0x19}
            \label{tab:test-dout-0x19} \\
            \hline
            \rowcolor{lightgray}
            \textbf{r-2'b00} & \textbf{r-1'b0} & \textbf{r-1'b0} & \textbf{r-4'b0000} \\ \hline
            Reserved & Bypass PLL loop divider default loopback & Bypass PLL input divider default loopback & Reserved \\ \hline
        \end{longtable}
\end{itemize}

\paragraph{HS RX Control of Lane 0}
\label{sec:test-code-0x44}
Controls the hsfreqrange.
\begin{itemize}
    \item \textbf{Test Code}: 0x44
    \item \textbf{Test Dout} 
        \begin{longtable}{ |p{2cm}|p{6cm}|p{6cm}|}
            \caption{Test Data of 0x44}
            \label{tab:test-data-0x44} \\
            \hline
            \rowcolor{lightgray}
            \textbf{w-1'b0} & \textbf{w-6'b000000} & \textbf{w-1'b0} \\ \hline
            Set 2 & HS operating frequency range selection (hsfreqrange) & Bypass HS-RX settle filter \\ \hline
        \end{longtable}

        \begin{itemize}
            \item Bits 7: Program selector
                \begin{itemize}
                    \item 0 - Set 2 is programmed
                    \item 1 - Invalid
                \end{itemize}
            \item Bit 6-1: HS frequency range selection. Refer to table \ref{tab:mipidsi-hsfreq} for detailed values.
            \item Bits 0: Bypass HS settle filter
        \end{itemize}
    \item \textbf{Test Dout}
        \begin{longtable}{ |p{2cm}|p{6cm}|p{6cm}|}
            \caption{Test Dout of 0x44}
            \label{tab:test-dout-0x44} \\
            \hline
            \rowcolor{lightgray}
            \textbf{r-1'b0} & \textbf{r-6'b000000} & \textbf{r-1'b0} \\ \hline
            1'b0 & HS operating frequency range selection loopback (hsfreqrange) & Bypass HS-RX settle filter loopback \\ \hline
        \end{longtable}
\end{itemize}


\subsection{Clocks and Display Resolution}
\label{sec:clocks-and-display-resolution}

The clock architecture of DSI is shown in figure \ref{figure:mipidsi-clk-arch}. All clocks are:
\begin{itemize}
    \item \textbf{PCLK}: the APB bus clock
    \item \textbf{ACLK}: the AXI bus clock
    \item \textbf{DSI\_DPICLK}: DPI interface clock
    \item \label{clk:dsi-dphy-cfg-clk}\textbf{DSI\_DPHY\_CFG\_CLK}: DPHY configuration clock
    \item \label{clk:dsi-dphy-pll-refclk}\textbf{DSI\_DPHY\_PLL\_REFCLK}: DPHY PLL input reference clock
    \item \textbf{TX\_HS\_CLK}: DPHY PLL output clock
    \item \textbf{TX\_BYTE\_CLK}: High-Speed (HS) transfer byte clock, TX\_HS\_CLK divided by 8.
    \item \textbf{TX\_ESC\_CLK}: Low-Power (LP) escape mode transfer clock, divided from TX\_BYTE\_CLK. The divider ratio can be configured through \hyperref[fielddesc:DSIHOSTTOCLKDIVISION]{DSI\_HOST\_TO\_CLK\_DIVISION}, and the allowed max-frequency is 20 MHz.
\end{itemize}

\begin{figure}[H]
    \centering
    \includegraphics[width=1.0\textwidth]{\modulefiles/figures/mipidsi-clk-arch.pdf}
    \caption{MIPI DSI clock architecture}
    \label{figure:mipidsi-clk-arch}
\end{figure}

DSI supports continuous and non-continuous clock-lane modes.
\begin{itemize}
    \item For continuous clock, set \hyperref[fielddesc:DSIHOSTAUTOCLKLANECTRL]{DSI\_HOST\_AUTO\_CLKLANE\_CTRL} to 1, set \hyperref[fielddesc:DSIHOSTPHYTXREQUESTCLKHS]{DSI\_HOST\_PHY\_TXREQUESTCLKHS} to 1
    \item For non-continuous clock, set \hyperref[fielddesc:DSIHOSTAUTOCLKLANECTRL]{DSI\_HOST\_AUTO\_CLKLANE\_CTRL} to 0, set \hyperref[fielddesc:DSIHOSTPHYTXREQUESTCLKHS]{DSI\_HOST\_PHY\_TXREQUESTCLKHS} to 1
\end{itemize}

In video mode transmission, the video timing parameters are shown in figure \ref{figure:mipidsi-video-timing}.
\begin{figure}[H]
    \centering
    \includegraphics[width=0.8\textwidth]{\modulefiles/figures/mipidsi-video-timing.pdf}
    \caption{Video Timing}
    \label{figure:mipidsi-video-timing}
\end{figure}

For DPI, the video timing and resolution can be configured through the following DSI Bridge register fields. After that, write 1 to \hyperref[fielddesc:DSIBRGDPICONFIGUPDATE]{DSI\_BRG\_DPI\_CONFIG\_UPDATE} to make the configuration valid.
\begin{itemize}
    \item \hyperref[fielddesc:DSIBRGHTOTAL]{DSI\_BRG\_HTOTAL}: The HTOTAL length measured in DSI\_DPICLK
    \item \hyperref[fielddesc:DSIBRGHDISP]{DSI\_BRG\_HDISP}: The HACT length measured in DSI\_DPICLK
    \item \hyperref[fielddesc:DSIBRGHBANK]{DSI\_BRG\_HBANK}: The HBP length measured in DSI\_DPICLK
    \item \hyperref[fielddesc:DSIBRGHSYNC]{DSI\_BRG\_HSYNC}: The HSA length measured in DSI\_DPICLK
    \item \hyperref[fielddesc:DSIBRGVTOTAL]{DSI\_BRG\_VTOTAL}: The VTOTAL length measured in lines
    \item \hyperref[fielddesc:DSIBRGVDISP]{DSI\_BRG\_VDISP}: The VACT length measured in lines
    \item \hyperref[fielddesc:DSIBRGVBANK]{DSI\_BRG\_VBANK}: The VBP length measured in lines
    \item \hyperref[fielddesc:DSIBRGVSYNC]{DSI\_BRG\_VSYNC}: The VSA length measured in lines
\end{itemize}


Similar configuration for video timing and resolution is required for DSI, and the related registers fields are listed below:
\begin{itemize}
    \item \hyperref[fielddesc:DSIHOSTVIDHSATIME]{DSI\_HOST\_VID\_HSA\_TIME}: The HSA length measured in TX\_BYTE\_CLK
    \item \hyperref[fielddesc:DSIHOSTVIDHBPTIME]{DSI\_HOST\_VID\_HBP\_TIME}: The HBP length measured in TX\_BYTE\_CLK
    \item \hyperref[fielddesc:DSIHOSTVIDHLINETIME]{DSI\_HOST\_VID\_HLINE\_TIME}: The HTOTAL length measured in TX\_BYTE\_CLK
    \item \hyperref[fielddesc:DSIHOSTVSALINES]{DSI\_HOST\_VSA\_LINES}: The VSA length measured in lines
    \item \hyperref[fielddesc:DSIHOSTVBPLINES]{DSI\_HOST\_VBP\_LINES}: The VBP length measured in lines
    \item \hyperref[fielddesc:DSIHOSTVFPLINES]{DSI\_HOST\_VFP\_LINES}: The VFP length measured in lines
    \item \hyperref[fielddesc:DSIHOSTVACTIVELINES]{DSI\_HOST\_VACTIVE\_LINES}: The VACT length measured in lines
\end{itemize}

The video timing configuration of DPI and DSI should match to avoid line buffer overflow. The actual timing of each line in DPI and DSI should be the same.

%%%%%%%%%%%%%%%%%%%%%
%%% Interrupts %%%
%%%%%%%%%%%%%%%%%%%%%

\section{Interrupts}
\chipname{}'s MIPI DSI can generate the following interrupt signals, which will be sent to the \href{mod:intmtrx}{Interrupt Matrix}.

\begin{itemize}
    \item dsi\_bridge\_interrupt
    \item dsi\_interrupt
\end{itemize}

The following interrupt sources can generate the dsi\_bridge\_interrupt interrupt signal:

\begin{longtable}[c]{ |p{5cm}|p{10cm}| }
    \caption{DSI Bridge interrupts}
    \label{tab:dsi-bridge-intr}\\\hline

    \rowcolor{lightgray}
    \textbf{Interrupt Source} & \textbf{Detail} \\\hline
    \endfirsthead


    \multicolumn{2}{c}{\bfseries \tablename\ \thetable{} -- ...Continued from the previous page}\\\hline
    \rowcolor{lightgray}
    \textbf{Interrupt Source} & \textbf{Detail} \\\hline
    \endhead

    \multicolumn{2}{r}{\textbf{Continued on the next page...}} \\
    \endfoot

    %\insertTableNotes
    \endlastfoot

    \hyperref[fielddesc:DSIBRGUNDERRUNINTRAW]{DSI\_BRG\_UNDERRUN\_INT} & Triggered when DSI Bridge buffer underruns \\\hline
        

\end{longtable}

The following interrupt sources can generate the dsi\_interrupt interrupt signal:



\begin{longtable}[c]{ |p{5cm}|p{5cm}|p{5cm}| }
    \caption{DSI Host Interrupts}
    \label{tab:dsi-host-interrupt}\\\hline

    \rowcolor{lightgray}
    \textbf{Interrupt} & \textbf{Cause of Interrupt} & \textbf{Recommended Handling}  \\\hline
    \endfirsthead


    \multicolumn{3}{c}{\bfseries \tablename\ \thetable{} -- ...Continued from the previous page}\\\hline
    \rowcolor{lightgray}
    \textbf{Interrupt} & \textbf{Cause of Interrupt} & \textbf{Recommended Handling}  \\\hline
    \endhead

    \multicolumn{3}{r}{\textbf{Continued on the next page...}} \\
    \endfoot

    %\insertTableNotes
    \endlastfoot

    \hyperref[fielddesc:DSIHOSTACKWITHERR0]{\makecell[l]{DSI\_HOST\_ACK\_WITH\_ERR\_0}} & SoT error from the Acknowledge error report & Transmit the packet again if the packet is critical \\\hline
    \hyperref[fielddesc:DSIHOSTACKWITHERR1]{\makecell[l]{DSI\_HOST\_ACK\_WITH\_ERR\_1}} & SoT Sync error from the Acknowledge error report & Transmit the packet again \\\hline
    \hyperref[fielddesc:DSIHOSTACKWITHERR2]{\makecell[l]{DSI\_HOST\_ACK\_WITH\_ERR\_2}} & EoT Sync error from the Acknowledge error report & Transmit the packet again \\\hline
    \hyperref[fielddesc:DSIHOSTACKWITHERR3]{\makecell[l]{DSI\_HOST\_ACK\_WITH\_ERR\_3}} & Escape Mode Entry Command error from the Acknowledge error report, the device does not recognize the Escape Mode Entry code & Check the device capability \\\hline
    \hyperref[fielddesc:DSIHOSTACKWITHERR4]{\makecell[l]{DSI\_HOST\_ACK\_WITH\_ERR\_4}} & Low-power Transmit Sync error from the Acknowledge error report & Transmit the packet again \\\hline
    \hyperref[fielddesc:DSIHOSTACKWITHERR5]{\makecell[l]{DSI\_HOST\_ACK\_WITH\_ERR\_5}} & Peripheral Timeout error from the Acknowledge error report & Check whether the Host timeout counter configuration is consistent with the Device specifications; Host timeout counter configuration should be longer than the Device. This error is automatically recovered, but with no guarantee that the transmission or reception is complete\\\hline
    \hyperref[fielddesc:DSIHOSTACKWITHERR6]{\makecell[l]{DSI\_HOST\_ACK\_WITH\_ERR\_6}} & False Control error from the Acknowledge error report & Reset the DSI Host and DPHY \\\hline
    \hyperref[fielddesc:DSIHOSTACKWITHERR7]{\makecell[l]{DSI\_HOST\_ACK\_WITH\_ERR\_7}} & Reserved (specific to device) from the Acknowledge error report & Check the Device capabilities and connectivity and transmit again \\\hline
    \hyperref[fielddesc:DSIHOSTACKWITHERR8]{\makecell[l]{DSI\_HOST\_ACK\_WITH\_ERR\_8}} & ECC error, single-bit (detected and corrected) from the Acknowledge error report & No action is required. Analyze the signal integrity or the noise conditions of the link if this error is recurrent \\\hline
    \hyperref[fielddesc:DSIHOSTACKWITHERR9]{\makecell[l]{DSI\_HOST\_ACK\_WITH\_ERR\_9}} & ECC error, multi-bit (detected, not corrected) from the Acknowledge error report & Transmit the packets again if the packets are critical. Analyze the signal integrity or the noise conditions of the link if this error is recurrent\\\hline
    \hyperref[fielddesc:DSIHOSTACKWITHERR10]{\makecell[l]{DSI\_HOST\_ACK\_WITH\_ERR\_10}} & Checksum error (long packet only) from the Acknowledge error report & Transmit the packets again if the packets are critical. Analyze the connectivity if this error is recurrent \\\hline
    \hyperref[fielddesc:DSIHOSTACKWITHERR11]{\makecell[l]{DSI\_HOST\_ACK\_WITH\_ERR\_11}} & Not recognized DSI data type from the Acknowledge error report. Some packets may not be supported & Check the Device capabilities \\\hline
    \hyperref[fielddesc:DSIHOSTACKWITHERR13]{\makecell[l]{DSI\_HOST\_ACK\_WITH\_ERR\_13}} & Invalid transmission length from the Acknowledge error report & Transmit the packets again. Analyze the connectivity if this error is recurrent \\\hline
    \hyperref[fielddesc:DSIHOSTACKWITHERR14]{\makecell[l]{DSI\_HOST\_ACK\_WITH\_ERR\_14}} & Reserved (specific to device) from the Acknowledge error report & Check the Device specification for possible reason and actions\\\hline
    \hyperref[fielddesc:DSIHOSTACKWITHERR15]{\makecell[l]{DSI\_HOST\_ACK\_WITH\_ERR\_15}} & DSI protocol violation from the Acknowledge error report & Check the display specification. With this error, the Device should have another read-back command which reports additional error information \\\hline
    \hyperref[fielddesc:DSIHOSTDPHYERRORS0]{\makecell[l]{DSI\_HOST\_DPHY\_ERRORS\_0}} & ErrEsc escape entry error from Lane 0. The Escape Entry Code is not recognized & Repeat the reception \\\hline
    \hyperref[fielddesc:DSIHOSTDPHYERRORS1]{\makecell[l]{DSI\_HOST\_DPHY\_ERRORS\_1}} & ErrSyncEsc low-power data transmission synchronization error from Lane 0 & Reset DSI Host and repeat the transmission \\\hline
    \hyperref[fielddesc:DSIHOSTDPHYERRORS2]{\makecell[l]{DSI\_HOST\_DPHY\_ERRORS\_2}} & Control error ErrControl from Lane 0 & Reset DSI Host and DPHY \\\hline
    \hyperref[fielddesc:DSIHOSTDPHYERRORS3]{\makecell[l]{DSI\_HOST\_DPHY\_ERRORS\_3}} & LP0 contention error ErrContentionLP0 from Lane 0 & Reset DSI Host and repeat the transmission \\\hline
    \hyperref[fielddesc:DSIHOSTDPHYERRORS4]{\makecell[l]{DSI\_HOST\_DPHY\_ERRORS\_4}} & LP1 contention error ErrContentionLP1 from Lane 0 & Reset DSI Host and repeat the transmission \\\hline
    \hyperref[fielddesc:DSIHOSTTOHSTX]{\makecell[l]{DSI\_HOST\_TO\_HS\_TX}} & High-speed transmission timeout error & Check the timeout configuration. Host timeout counter configuration should be longer than the Device \\\hline
    \hyperref[fielddesc:DSIHOSTTOLPRX]{\makecell[l]{DSI\_HOST\_TO\_LP\_RX}} & Low-power reception timeout error & Check the timeout configuration. Host timeout counter configuration should be longer than the Device \\\hline
    \hyperref[fielddesc:DSIHOSTECCSINGLEERR]{\makecell[l]{DSI\_HOST\_ECC\_SINGLE\_ERR}} & ECC single error has been detected and corrected in a received packet & No action is required \\\hline
    \hyperref[fielddesc:DSIHOSTECCMILTIERR]{\makecell[l]{DSI\_HOST\_ECC\_MILTI\_ERR}} & ECC multiple error has been detected in a received packet & Repeat the transmission \\\hline
    \hyperref[fielddesc:DSIHOSTCRCERR]{\makecell[l]{DSI\_HOST\_CRC\_ERR}} & CRC error has been detected in the received packet payload & Reset DSI Host and repeat the transmission \\\hline
    \hyperref[fielddesc:DSIHOSTPKTSIZEERR]{\makecell[l]{DSI\_HOST\_PKT\_SIZE\_ERR}} & Packet size error has been detected during the packet reception & Reset DSI Host and repeat the transmission \\\hline
    \hyperref[fielddesc:DSIHOSTEOPTERR]{\makecell[l]{DSI\_HOST\_EOPT\_ERR}} & EoTp packet has not been received at the end of the incoming peripheral transmission & This error is not critical for data integrity of the received packet. Check if the Device supports EoTp packets \\\hline
    \hyperref[fielddesc:DSIHOSTDPIPLDWRERR]{\makecell[l]{DSI\_HOST\_DPI\_PLD\_WR\_ERR}} & During a DPI pixel line storage, the payload FIFO becomes full and the data stored is corrupted & Check Video Mode configuration, and reset DSI Host and re-initiate the video transmission \\\hline
    \hyperref[fielddesc:DSIHOSTGENCMDWRERR]{\makecell[l]{DSI\_HOST\_GEN\_CMD\_WR\_ERR}} & The system tried to write a command through the Generic interface and the FIFO is full. Therefore, the command is not written & Reset DSI Host and repeat the transmission \\\hline
    \hyperref[fielddesc:DSIHOSTGENPLDWRERR]{\makecell[l]{DSI\_HOST\_GEN\_PLD\_WR\_ERR}} & The system tried to write a payload data through the Generic interface and the FIFO is full. Therefore, the payload is not written & Reset DSI Host and repeat the transmission \\\hline
    \hyperref[fielddesc:DSIHOSTGENPLDSENDERR]{\makecell[l]{DSI\_HOST\_GEN\_PLD\_SEND\\\_ERR}} & During a Generic interface packet build, the payload FIFO becomes empty and corrupt data is sent & Reset DSI Host and repeat the transmission \\\hline
    \hyperref[fielddesc:DSIHOSTGENPLDRDERR]{\makecell[l]{DSI\_HOST\_GEN\_PLD\_RD\_ERR}} & During a DCS read data, the payload FIFO becomes empty and the data sent to the interface is corrupted & Reset DSI Host and repeat the transmission \\\hline
    \hyperref[fielddesc:DSIHOSTGENPLDRECEVERR]{\makecell[l]{DSI\_HOST\_GEN\_PLD\_RECEV\\\_ERR}} & During a Generic interface packet read back, the payload FIFO becomes full and the received data is corrupted & Reset DSI Host and repeat the transmission \\\hline
    \hyperref[fielddesc:DSIHOSTDPIBUFFPLDUNDER]{\makecell[l]{DSI\_HOST\_DPI\_BUFF\_PLD\\\_UNDER}} & An underflow has occurred when reading payload to build DSI packet for video mode & Reset DSI Host and repeat the transmission \\\hline

\end{longtable}

%%%%%%%%%%%%%%%%%%%%%%%%%%%%%%
%%% Programming Procedures %%%
%%%%%%%%%%%%%%%%%%%%%%%%%%%%%%

\section{Programming Procedures}
\label{sec:mipidsi-prog-procedures}
\subsection{DSI Bridge Configuration}
\label{sec:dsi-bridge-configuration}
\begin{itemize}
    \item Configure DSI Bridge clock and release reset.
    \item Configure input and output color mode. Refer to \hyperref[sec:dsi-bridge]{DSI Bridge} for details.
    \item Configure video timing. Refer to \hyperref[sec:clocks-and-display-resolution]{Clocks and Display Resolution} for details.
    \item Write 0 to \hyperref[fielddesc:DSIBRGDSIDMAFLOWCONTROLLER]{DSI\_BRG\_DSI\_DMA\_FLOW\_CONTROLLER} to set VDMA as flow controller.
    \item Configure \hyperref[fielddesc:DSIBRGDMABURSTLEN]{DSI\_BRG\_DMA\_BURST\_LEN} to set the DMA burst length to match VDMA configuration.
    \item Configure \hyperref[fielddesc:DSIBRGDSIRAWBUFALMOSTEMPTYTHRD]{DSI\_BRG\_DSI\_RAW\_BUF\_ALMOST\_EMPTY\_THRD} to $1024 - \hyperref[fielddesc:DSIBRGDMABURSTLEN]{DSI\_BRG\_DMA\_BURST\_LEN}$ for better performance.
    \item Write 1 to \hyperref[fielddesc:DSIBRGDPIEN]{DSI\_BRG\_DPI\_EN} to enable DPI.
    \item Write 1 to \hyperref[fielddesc:DSIBRGDSIEN]{DSI\_BRG\_DSI\_EN} to enable DSI Bridge.
    \item Configure VDMA and set VDMA as flow controller. Refer to Chapter \ref{mod:vdma} \textit{\nameref{mod:vdma}} for details.
\end{itemize}
\label{sec:mipidsi-dsi-bridge-programming-proc}

\subsection{D-PHY Configuration}
\label{sec:mipidsi-dsi-dphy-programming-proc}
\begin{itemize}
    \item Set \hyperref[fielddesc:DSIHOSTPHYSHUTDOWNZ]{DSI\_HOST\_PHY\_SHUTDOWNZ} and \hyperref[fielddesc:DSIHOSTPHYRSTZ]{DSI\_HOST\_PHY\_RSTZ} to 0, set \hyperref[fielddesc:DSIHOSTPHYTESTCLR]{DSI\_HOST\_PHY\_TESTCLR} to 1.
    \item Apply an appropriate \hyperref[clk:dsi-dphy-cfg-clk]{DSI\_DPHY\_CFG\_CLK} and \hyperref[clk:dsi-dphy-pll-refclk]{DSI\_DPHY\_PLL\_REFCLK} (see Table~\ref{tab:resclk-periph-clks-hp2} in Chapter~\ref{mod:resclk} \textit{\nameref{mod:resclk}} for clock sources and frequency requirements).
    \item Configure the hsfreqrange through \hyperref[sec:test-code-0x44]{test code 0x44}. Refer to table \ref{tab:mipidsi-hsfreq} for detailed value.
    \item Configure PLL parameters. Refer to \hyperref[sec:pll-parameters]{PLL Parameters} for details.
    \item Configure \hyperref[fielddesc:DSIHOSTTOCLKDIVISION]{DSI\_HOST\_TO\_CLK\_DIVISION} to apply an appropriate TX\_ESC\_CLK divider ratio.
    \item Set \hyperref[fielddesc:DSIHOSTPHYSHUTDOWNZ]{DSI\_HOST\_PHY\_SHUTDOWNZ} and \hyperref[fielddesc:DSIHOSTPHYRSTZ]{DSI\_HOST\_PHY\_RSTZ} to 1.
    \item Wait until \hyperref[fielddesc:DSIHOSTPHYSTOPSTATE0LANE]{DSI\_HOST\_PHY\_STOPSTATE0LANE} and \hyperref[fielddesc:DSIHOSTPHYSTOPSTATE1LANE]{DSI\_HOST\_PHY\_STOPSTATE1LANE} are asserted.
\end{itemize}

\subsection{MIPI DSI Common Configuration}
\label{sec:mipidsi-dsi-common-programming-proc}
\begin{itemize}
    \item Release global reset.
    \item Configure \hyperref[fielddesc:DSIHOSTNLANES]{DSI\_HOST\_N\_LANES} to define the number of lanes to be used.
    \item Configure video timing settings. Refer to \hyperref[sec:clocks-and-display-resolution]{Clocks and Display Resolution} for details.
    \item Initialize the DPHY through \hyperref[sec:mipidsi-dsi-dphy-programming-proc]{D-PHY Configuration}.
    \item Write \hyperref[fielddesc:DSIHOSTSHUTDOWNZ]{DSI\_HOST\_SHUTDOWNZ} to 1 to release DSI Host from shutdown.
\end{itemize}

\subsection{System memory data output flow configuration}
\begin{itemize}
    \item Configure DSI Bridge through \hyperref[sec:dsi-bridge-configuration]{DSI Bridge Configuration}.
    \item Configure DSI Host through \hyperref[sec:mipidsi-dsi-common-programming-proc]{MIPI DSI Common Configuration}.
\end{itemize}

\subsection{Video pattern output flow configuration}
\begin{itemize}
    \item Configure pattern through register field \hyperref[fielddesc:DSIHOSTVPGORIENTATION]{DSI\_HOST\_VPG\_ORIENTATION}, \hyperref[fielddesc:DSIHOSTVPGMODE]{DSI\_HOST\_VPG\_MODE}, and \hyperref[fielddesc:DSIHOSTVPGEN]{DSI\_HOST\_VPG\_EN}. Refer to \hyperref[sec:video-pattern-generator]{Video Pattern Generator} for details.
    \item Configure DSI Host through \hyperref[sec:mipidsi-dsi-common-programming-proc]{MIPI DSI Common Configuration}.
\end{itemize}



\subsection{Using DCS command to operate Display register}
\begin{itemize}
    \item Write to Display
        \begin{itemize}
            \item Write packet payload to \hyperref[regdesc:DSIHOSTGENPLDDATAREG]{DSI\_HOST\_GEN\_PLD\_DATA\_REG}.
            \item Write packet header to \hyperref[regdesc:DSIHOSTGENHDRREG]{DSI\_HOST\_GEN\_HDR\_REG}.
        \end{itemize}
    \item Read from Display
        \begin{itemize}
            \item Write packet header to \hyperref[regdesc:DSIHOSTGENHDRREG]{DSI\_HOST\_GEN\_HDR\_REG}.
            \item Read response data from \hyperref[regdesc:DSIHOSTGENPLDDATAREG]{DSI\_HOST\_GEN\_PLD\_DATA\_REG}.
        \end{itemize}
\end{itemize}

% >>> If your module has no registers, delete sample text
%       for Sections Register Summary and Registers below <<<

%%%%%%%%%%%%%%%%%%%%%%%%
%%% Register Summary %%%
%%%%%%%%%%%%%%%%%%%%%%%%

% >>> Update the sample text below and insert
%     the info on register summary into the subfile <<<

\newpage
\hypertarget{mipidsi-reg-summ}{}
\section{Register Summary}
\label{sec:mipidsi-reg-summ}


\subsection{MIPI DSI Bridge Register Summary}

The addresses in this section are relative to MIPI DSI BRIDGE base address provided in Table \ref{tab:sysmem-base-address} in Chapter \ref{mod:sysmem} \textit{\nameref{mod:sysmem}}.

The abbreviations given in Column \textbf{Access} are explained in Section \hyperref[glossary-access-types]{\textit{Access Types for Registers}}.

\subfile{\modulefiles/42-MIPIDSI-BRG-reg-summ__EN}

\subsection{MIPI DSI Host Register Summary}

The addresses in this section are relative to MIPI DSI HOST base address provided in Table \ref{tab:sysmem-base-address} in Chapter \ref{mod:sysmem} \textit{\nameref{mod:sysmem}}.

The abbreviations given in Column \textbf{Access} are explained in Section \hyperref[glossary-access-types]{\textit{Access Types for Registers}}.

\subfile{\modulefiles/42-MIPIDSI-reg-summ__EN}



%%%%%%%%%%%%%%%%%
%%% Registers %%%
%%%%%%%%%%%%%%%%%

% >>> Update the sample text below and insert
%      the info on registers into the subfile <<<

\newpage
\section{Registers}

\subsection{MIPI DSI Bridge Register}
The addresses in this section are relative to MIPI DSI BRIDGE base address provided in Table \ref{tab:sysmem-base-address} in Chapter \ref{mod:sysmem} \textit{\nameref{mod:sysmem}}.

For how to program reserved fields, please refer to Section \hyperref[programming-reserved-field]{Programming Reserved Register Field}.

\subfile{\modulefiles/42-MIPIDSI-BRG-reg__EN}

\subsection{MIPI DSI Host Register}
The addresses in this section are relative to MIPI DSI HOST base address provided in Table \ref{tab:sysmem-base-address} in Chapter \ref{mod:sysmem} \textit{\nameref{mod:sysmem}}.

For how to program reserved fields, please refer to Section \hyperref[programming-reserved-field]{Programming Reserved Register Field}.

\subfile{\modulefiles/42-MIPIDSI-reg__EN}

\end{document}
